The equality set $\{\mathbf{x}:g_{jl}(\mathbf{x};\boldsymbol{\theta})=0\}$ is not necessarily a boundary of
the classification partition.  A third component can dominate both $j$ and
$l$ on part of that set.  For $j<l$, define the relatively open active face
\begin{equation}
 \mathcal F_{jl}(\boldsymbol{\theta})
 =\left\{\mathbf{x}:q_j(\mathbf{x};\boldsymbol{\theta})=q_l(\mathbf{x};\boldsymbol{\theta})>
 q_m(\mathbf{x};\boldsymbol{\theta})\ \text{for every }m\notin\{j,l\}\right\}.
 \label{eq:active-face}
\end{equation}
Only $\mathcal F_{jl}(\boldsymbol{\theta})$ separates cells $j$ and $l$ locally.  Its
closure can meet other faces at multiple-tie junctions.  We write
$\mathcal F_{jl}^*=\mathcal F_{jl}(\boldsymbol{\theta}^*)$.

For a smooth scalar map $g$, the notation $\nabla_{\mathbf{x}}g$ and $\nabla_{\boldsymbol{\theta}}g$ denotes
gradients with respect to $\mathbf{x}$ and $\boldsymbol{\theta}$, respectively; $\nabla_{\boldsymbol{\theta}}^2q_j$ denotes
the Hessian with respect to $\boldsymbol{\theta}$.  Let $\mathcal H^{p-1}$ denote
$(p-1)$-dimensional Hausdorff measure.

\begin{assumption}[Regular active-boundary geometry]
\label{ass:boundary-geometry}
There is an open neighbourhood $U$ of the interior point $\boldsymbol{\theta}^*$ such that:
\begin{enumerate}
 \item $P_0$ has a Lebesgue density $p_0$ that is continuous on a
 neighbourhood of the closures of the active faces at $\boldsymbol{\theta}^*$.
 \item Each $q_j(\mathbf{x};\boldsymbol{\theta})$ is twice continuously differentiable in $\boldsymbol{\theta}$
 and continuously differentiable in $\mathbf{x}$ on the relevant neighbourhood.
 \item For every nonempty $\mathcal F_{jl}^*$,
 $\nabla_{\mathbf{x}}g_{jl}(\mathbf{x};\boldsymbol{\theta}^*)\ne0$ on its closure except possibly on a set of
 $\mathcal H^{p-1}$-measure zero.
 \item The set on which three or more component scores attain the common
 maximum has $\mathcal H^{p-1}$-measure zero.
 \item The derivatives can be passed through the volume and surface
 integrals below.  In particular, the interior Hessian terms are
 $P_0$-integrable and
 \[
  \int_{\mathcal F_{jl}^*}
  \frac{p_0(\mathbf{x})\|\nabla_{\boldsymbol{\theta}}g_{jl}(\mathbf{x};\boldsymbol{\theta}^*)\|^2}
       {\|\nabla_{\mathbf{x}}g_{jl}(\mathbf{x};\boldsymbol{\theta}^*)\|}
  \,d\mathcal H^{p-1}(\mathbf{x})<\infty
 \]
 for every $j<l$.
\end{enumerate}
\end{assumption}

Assumption~\ref{ass:boundary-geometry}(4) is a geometric condition, not merely
a zero-$P_0$-probability condition.  Every regular hypersurface has
$P_0$-probability zero under a Lebesgue density, while it can have positive
$\mathcal H^{p-1}$ measure and contribute to curvature.

At points with a unique winning component, define
\begin{equation}
 \boldsymbol{\psi}(\mathbf{x};\boldsymbol{\theta})
 =\sum_{j=1}^k
 \ind\{a_{\boldsymbol{\theta}}(\mathbf{x})=j\}\nabla_{\boldsymbol{\theta}}q_j(\mathbf{x};\boldsymbol{\theta}).
 \label{eq:classification-score}
\end{equation}
On ties the right-hand side uses the deterministic convention from
\eqref{eq:classification-cells}.  Its value on those points is immaterial under
Assumption~\ref{ass:boundary-geometry}.

\begin{theorem}[Population derivatives and active-boundary curvature]
\label{thm:population-curvature}
Suppose Assumption~\ref{ass:boundary-geometry} holds and the indicated
first-order derivatives admit integrable envelopes.  Then $M$ is differentiable
at $\boldsymbol{\theta}^*$ with
\begin{equation}
 \nabla_{\boldsymbol{\theta}}M(\boldsymbol{\theta}^*)=P_0\boldsymbol{\psi}(\mathbf{X};\boldsymbol{\theta}^*).
 \label{eq:population-gradient}
\end{equation}
It is twice differentiable at $\boldsymbol{\theta}^*$, and
\begin{align}
 \mathbf{H}:=\nabla_{\boldsymbol{\theta}}^2M(\boldsymbol{\theta}^*)
 &=P_0\left[
   \sum_{j=1}^k\ind\{a_{\boldsymbol{\theta}^*}(\mathbf{X})=j\}
   \nabla_{\boldsymbol{\theta}}^2q_j(\mathbf{X};\boldsymbol{\theta}^*)
  \right]+\mathbf{B},                                      \label{eq:hessian-decomposition}\\
 \mathbf{B}
 &=\sum_{1\leq j<l\leq k}
   \int_{\mathcal F_{jl}^*}
   \frac{p_0(\mathbf{x})\nabla_{\boldsymbol{\theta}}g_{jl}(\mathbf{x};\boldsymbol{\theta}^*)
        \nabla_{\boldsymbol{\theta}}g_{jl}(\mathbf{x};\boldsymbol{\theta}^*)^\mathsf T}
        {\|\nabla_{\mathbf{x}}g_{jl}(\mathbf{x};\boldsymbol{\theta}^*)\|}
   \,d\mathcal H^{p-1}(\mathbf{x}).                                \label{eq:boundary-curvature}
\end{align}
Consequently, if
\begin{equation}
 \mathbf{I}_c=-P_0\left[
   \sum_{j=1}^k\ind\{a_{\boldsymbol{\theta}^*}(\mathbf{X})=j\}
   \nabla_{\boldsymbol{\theta}}^2q_j(\mathbf{X};\boldsymbol{\theta}^*)
  \right],
 \label{eq:fixed-class-information}
\end{equation}
then the positive population curvature is
\begin{equation}
 \mathbf{A}=-\mathbf{H}=\mathbf{I}_c-\mathbf{B}.
 \label{eq:positive-curvature}
\end{equation}
In particular, $\mathbf{B}$ is positive semidefinite and parameter-dependent
classification weakens the positive curvature relative to fixed-label
analysis.
\end{theorem}

\begin{proof}
Away from ties, the maximizing index is locally constant in $\boldsymbol{\theta}$, so the
first derivative of $m_{\boldsymbol{\theta}}(\mathbf{x})$ equals the derivative of the winning score.
The envelope condition and the fact that regular ties have $P_0$-probability
zero give \eqref{eq:population-gradient} by dominated differentiation.

For the second derivative, first consider a neighbourhood containing only an
active face between $j$ and $l$.  Locally,
\[
 m_{\boldsymbol{\theta}}(\mathbf{x})=q_l(\mathbf{x};\boldsymbol{\theta})+\{g_{jl}(\mathbf{x};\boldsymbol{\theta})\}_+.
\]
In the distributional derivative with respect to $\boldsymbol{\theta}$, differentiating
the positive-part term a second time produces the ordinary within-region
Hessian and
\[
 \delta\{g_{jl}(\mathbf{x};\boldsymbol{\theta}^*)\}
 \nabla_{\boldsymbol{\theta}}g_{jl}(\mathbf{x};\boldsymbol{\theta}^*)\nabla_{\boldsymbol{\theta}}g_{jl}(\mathbf{x};\boldsymbol{\theta}^*)^\mathsf T.
\]
The coarea formula \citep{evans2015measure} converts its integral with respect to
$p_0(\mathbf{x})\,dx$ into the corresponding term in
\eqref{eq:boundary-curvature}.  A partition of unity over the unique-winner
regions and the relatively open active faces combines these local identities.
Inactive portions of pairwise equality sets do not enter because the maximum
is locally attained by another component there.  Multiple-tie junctions make
no contribution by Assumption~\ref{ass:boundary-geometry}(4).  Summing over
$j<l$ proves \eqref{eq:hessian-decomposition}; the remaining statements follow
from the definitions.
\end{proof}

\begin{remark}[Relation to $k$-means]
\label{rem:kmeans}
For $q_j(\mathbf{x};\boldsymbol{\theta})=-\|\mathbf{x}-
\boldsymbol{\mu}_j\|^2$, the active faces are Voronoi boundaries.  Pollard's
$k$-means boundary geometry is therefore recovered as a special
hard-classification case \citep{pollard1982clt}.  The result extends it to
weighted log-density comparisons, heterogeneous families, and misspecification.
\end{remark}

\begin{corollary}[Negligible-boundary limit]
\label{cor:negligible-boundary}
If $p_0=0$ $\mathcal H^{p-1}$-almost everywhere on every active face, then
$\mathbf{B}=0$ and fixed-classification curvature is correct to first order.
More generally, the size of the correction is governed by probability density
on the active faces, the sensitivity $\nabla_{\boldsymbol{\theta}}g_{jl}$, and the spatial
transversality $\|\nabla_{\mathbf{x}}g_{jl}\|$.
\end{corollary}
